\documentclass[12pt]{article}

\title{Homework week 3}
\author{
	Evan M Drake \\
	\texttt{}
}


\usepackage{snow-preamble}

%%%%%%%%%%%%%%%%%%%%%%%%%%%%%%%%%%%%
\begin{document}
	\maketitle
	% \tableofcontents

	% \newpage

	%% body start

	\section{Question 1}
		For the following questions give both the output and its type.
		For example if the question is $a + b = c$, $c$ is the output and its type is assumed to be number or integer.
		The questions are given in MATLAB code,
			so they should run,
			if they do not run then the error should be given as the answer.

		\begin{enumerate}
			\item $5^5$
			\item $16 + 30$
			\item $32/4$
			\item $25/2 * 6$
		\end{enumerate}

		\setcounter{equation}{0}
	\section{Question 2}
		Write MATLAB code for the following equations.

		\begin{align}
			5 + 5 \\
			\sqrt{4} \\
			3^1.2 \\
			\tan{60}
		\end{align}

		\setcounter{equation}{0}
	\section{Question 3}
		Store a variable $ftemp$ giving it an initial value of 60.
		We will assume $ftemp$ is the degrees of a room in Fahrenheit.
		Create a fxn to convert numbers from Fahrenheit to Celsius.
		Invoke your conversion fxn on four different variables including $ftemp$.
		Turn in your code as the answer.

	\section{Question 4}
		Write MATLAB expressions to generate a random number in the following ranges.
		\begin{align*}
			1 \rightarrow 10 \\
			1 \rightarrow 100 \\
			0 \rightarrow 5 \\
			-100 \rightarrow 100
		\end{align*}

	\section{Question 5}
		Create the following matrix below.
		\begin{equation}
			M = \mathcal{R}^{10x10}
		\end{equation}

		\setcounter{equation}{0}
	\section{Question 6}
		For the matrix $M$ defined in question 8,
			do the following.
		\begin{enumerate}
			\item return the value that is at index $M_{1,1}$
			\item return the value that is at index $M_{10,10}$
			\item return the second row
			\item return the third column
		\end{enumerate}
	
	\section{Question 7}
		Write the pseudocode for a function that would theoretically perform matrix multiplication.
		
	%% body finish

	\nocite{*}

	\newpage % bib
	\backmatter
	\addcontentsline{toc}{section}{References}
	\printbibliography

	\newpage % index
	\printindex

	\newpage % appendices
	\appendix

\end{document}

